\section{结论}
\label{sec:conclusion}

本文围绕检索触发、查询构建、证据评价和停止决策梳理了自适应 RAG。输入路由、生成中触发、反思标记和迭代检索提供了不同控制粒度，搜索策略学习、过程监督与开放网络研究又扩展了行动空间。检索必要性随模型、语料和任务要求变化，证据充分性也不能由文档数量或槽位填充比例单独确定。

停止控制中的内部信念、证据覆盖和答案效用各有观测对象与评价边界，尤其部分效用依赖答案探针和质量指标。评价自适应 RAG 时，需要同时说明任务目标、充分性定义和调用成本，并结合跨域测试与轨迹干预判断策略学到的是可迁移的信息获取能力，还是特定基准上的行为偏好。
